\documentclass[12pt,a4paper]{article}

% ---------------------------------------------------
% GRUNDSPRACHE & SCHRIFT
% ---------------------------------------------------
\usepackage[ngerman]{babel}
\usepackage{fontspec}
\setmainfont{Times New Roman} % Times New Roman

% ---------------------------------------------------
% SEITENLAYOUT
% ---------------------------------------------------
\usepackage{geometry}
\geometry{
    left=4cm,
    right=2cm,
    top=4cm,
    bottom=2cm
}

% ---------------------------------------------------
% FORMATIERUNG
% ---------------------------------------------------
\usepackage{setspace}
\onehalfspacing
\setlength{\parskip}{6pt}
\setlength{\parindent}{0pt}

\usepackage{ragged2e}
\justifying
\hyphenation{}

\usepackage{titlesec}
\titlespacing*{\section}{0pt}{12pt}{6pt}
\titlespacing*{\subsection}{0pt}{12pt}{6pt}
\usepackage{etoolbox}
\pretocmd{\section}{\clearpage}{}{}


% ---------------------------------------------------
% GRAFIKEN, TABELLEN, MATHE
% ---------------------------------------------------
\usepackage{graphicx}
\usepackage{booktabs}
\usepackage{amsmath}
\usepackage{caption}

% ---------------------------------------------------
% CODEVERZEICHNIS
% ---------------------------------------------------
\usepackage{listings}
\renewcommand{\lstlistlistingname}{Codeverzeichnis}
\lstset{
    basicstyle=\ttfamily\small,
    breaklines=true,
    frame=single,
    captionpos=b
}

% ---------------------------------------------------
% ZITATE, LINKS, VERZEICHNISSE
% ---------------------------------------------------
\usepackage{csquotes}
\usepackage[hidelinks]{hyperref}
\usepackage[nottoc]{tocbibind} % << WICHTIG: Inhaltsverzeichnis NICHT im ToC

% ---------------------------------------------------
% KOPF-/FUSSZEILEN
% ---------------------------------------------------
\usepackage{fancyhdr}
\pagestyle{fancy}
\fancyhf{}
\renewcommand{\headrulewidth}{0pt}
\fancyhead[C]{\thepage}
\setlength{\headheight}{15pt}
\setlength{\headsep}{20pt}

\usepackage[
backend=biber,
style=authoryear,
citestyle=authoryear,
sorting=nyt,
maxcitenames=2,
maxbibnames=99,
giveninits=true,
uniquename=false,
uniquelist=false,
doi=false,
isbn=false,
url=true
]{biblatex}
\addbibresource{literatur.bib}

% ---------------------------------------------------
% DOKUMENTBEGINN
% ---------------------------------------------------
\begin{document}

% ---------------------------------------------------
% TITELBLATT
% ---------------------------------------------------
\begin{titlepage}
\begin{center}
  \IfFileExists{FOM_Logo01_2012_rgb.jpg}{%
    \includegraphics[width=4cm]{FOM_Logo01_2012_rgb.jpg} \\[0.5cm]
  }{%
    {\Large \textbf{FOM}}\\[0.5cm]
  }
    \textbf{FOM Hochschule für Oekonomie \& Management} \\
    Hochschulzentrum Bonn \\[1.0cm]

    {\Large Berufsbegleitender Studiengang Wirtschaftsinformatik} \\[0.2cm]
    zum \\[0.2cm]
    Bachelor of Science (B.Sc.) \\[0.2cm]
    6. Semester \\[1.5cm]

    Seminararbeit im Modul Anwendungsprojekt zum Thema \\[1.5cm]

    {\Large \textbf{„Analyse der Auswirkung geopolitischer Konflikte auf \\ Aktienindizes“}} \\[1.5cm]
\end{center}

\vfill

\begin{flushleft}
\begin{tabular}{@{}ll} 
Dozent: & Prof. Dr. Julian Schröter \\
Name: & Lorenz Chylka, Leonard Reglin, \\
      & Lukas Nießen, Johannes Solbach \\
Matrikelnummer: & 736986, 740270, 728085, 732936 \\
Abgabedatum: & 15.06.2026 \\
Wörter Hausarbeit: & XXX \\
Wörter Code: & XXX \\
\end{tabular}
\end{flushleft}
\end{titlepage}

% ---------------------------------------------------
% VERZEICHNISSE (RÖMISCH)
% ---------------------------------------------------
\pagenumbering{Roman}
\setcounter{page}{2}
\tableofcontents
\newpage

\listoffigures
\newpage

\listoftables
\newpage

\lstlistoflistings
\newpage

\section*{Abkürzungsverzeichnis}
\addcontentsline{toc}{section}{Abkürzungsverzeichnis}
\begin{tabular}{@{}ll}
CRISP-DM & Cross Industry Standard Process for Data Mining \\
DBSCAN & Density-Based Spatial Clustering of Applications with Noise \\
KNN & K-Nearest Neighbor Algorithmus \\
ML & Machine Learning \\
GPR & Geopolitical Risk Index \\
\end{tabular}
\newpage
% ---------------------------------------------------
% LITERATUR-VERZEICHNIS
% ---------------------------------------------------
\newpage
\printbibliography[heading=bibintoc, title={Literaturverzeichnis}]
% ---------------------------------------------------
% KI-VERZEICHNIS
% ---------------------------------------------------
\newpage
\section*{KI-Hilfsmittelverzeichnis}
\addcontentsline{toc}{section}{KI-Hilfsmittelverzeichnis}

\begin{tabular}{@{} l l p{9.5cm} @{}}
\toprule
\textbf{Nr.} & \textbf{KI-System} & \textbf{Prompt / Einsatzzweck} \\ \midrule
1 & Google Gemini & „Hallo Gemini, ich brauche für mein Anwendungsprojekt eine wissenschaftliche Einleitung. Wir untersuchen, wie geopolitische Krisen Aktienkurse beeinflussen. Kannst du einen Text schreiben, der erklärt, warum das wichtig ist (Stichwort Volatilität und Unsicherheit) und warum wir genau Machine Learning dafür benutzen? Das Ganze soll ca. 250 Wörter haben.“ \\
2 & Google Gemini & „Hilf mir mal bitte beim 'Business Understanding' für meine Hausarbeit. Wir nutzen das CRISP-DM Modell und ich muss erklären, was das Ziel unserer Analyse ist. Kannst du mir da Hypothesen formulieren, wie Märkte auf Konflikte reagieren, und den wirtschaftlichen Nutzen für Investoren herausstellen?“ \\
3 & Google Gemini & „Ich muss das Kapitel 'Data Understanding' schreiben. Ich habe Aktienmarktdaten von Kaggle und Konflikt-Events von der UCDP. Kannst du beschreiben, was diese Quellen beinhalten und wie wir sie zusammenführen? Integriere bitte auch die Datenschemata-Tabellen, damit man sieht, welche Spalten wir genau nutzen.“ \\ 
4 & ChatGPT-5.4 & erläutere die data preparation, gehe auch auf die dafür genutzen skripte auf. erläutere, warum dies so gemacht wurde. erläutere dies auf einem wissenschaftlichen level und packe dies in das latex dokument bei 2.3 Data Preparation \\

git checkout NEU_k_means
error: pathspec 'NEU_k_means' did not match any file(s) known to git\bottomrule
\end{tabular}


% ---------------------------------------------------
% TEXTTEIL
% ---------------------------------------------------
\newpage
\pagenumbering{arabic}
\setcounter{page}{1}

\section{Einleitung}
Die globale Finanzwelt wird in zunehmendem Maße von unvorhersehbaren geopolitischen Ereignissen geprägt, die oft als „Black Swan“-Ereignisse die Stabilität internationaler Märkte nachhaltig erschüttern können. Geopolitische Risiken, sei es in Form von eskalierenden bewaffneten Konflikten, grenzüberschreitenden terroristischen Anschlägen oder tiefgreifenden diplomatischen Spannungen, induzieren eine signifikante Unsicherheit unter den globalen Marktteilnehmern. Diese Unsicherheit spiegelt sich unmittelbar in einer erhöhten Marktvolatilität, massiven Kapitalumschichtungen und oft abrupten Kurskorrekturen in führenden Leitindizes wider \parencite{caldara2022measuring}. In einer Zeit, in der Informationen in Echtzeit verarbeitet werden, ist die Fähigkeit, solche Schocks quantitativ zu erfassen, für institutionelle Investoren von existenzieller Bedeutung.

Klassische ökonometrische Modelle stoßen bei der Analyse dieser komplexen Zusammenhänge jedoch häufig an ihre methodischen Grenzen. An dieser Stelle bietet das Machine Learning (ML) innovative Lösungsansätze, um verborgene Strukturen und Korrelationen in hochdimensionalen Datenräumen zu identifizieren, die für herkömmliche statistische Tests unsichtbar bleiben. Im Zentrum dieses Anwendungsprojekts steht daher die Entwicklung eines ML-gestützten Frameworks zur Analyse geopolitischer Marktreaktionen.

Dabei kommen komplementäre Algorithmen zum Einsatz: Die lineare Regression dient der Bestimmung grundlegender Sensitivitäten, während DBSCAN zur Identifikation von Marktphasen und Anomalien genutzt wird. Ergänzend ermöglicht der k-Nearest-Neighbor-Algorithmus (KNN) einen fundierten historischen Vergleich, indem er aktuelle Marktszenarien mit ähnlichen Mustern aus der Vergangenheit in Beziehung setzt. Das Ziel der Arbeit ist die Schaffung einer transparenten Datenpipeline, die die Resilienz globaler Indizes gegenüber geopolitischen Schocks objektiv bewertbar macht und somit einen Beitrag zum modernen Risikomanagement leistet. (in Anlehnung an Google Gemini)

% ===================================================
\section{Projektbeschreibung nach CRISP-DM}

\subsection{Business Understanding}
Die Phase des Business Understanding bildet das Fundament des Projekts nach dem CRISP-DM-Standardmodell \parencite{wirth2000crisp}. In diesem ersten Schritt werden die Projektziele aus einer fachlichen Perspektive definiert. Das primäre Ziel besteht darin, die Sensitivität globaler Finanzmärkte gegenüber geopolitischen Schocks quantitativ zu bewerten. 

Die zentrale Forschungsfrage lautet: Inwiefern lassen sich durch die Kombination historischer Konfliktdaten (UCDP) und Marktdaten signifikante Muster der Marktreaktion identifizieren?
Daraus leiten sich folgende Hypothesen ab:
\begin{itemize}
    \item H1: Geopolitische Krisen führen zu signifikanten, aber kurzfristigen Kursrückgängen in Leitindizes wie dem S\&P 500 oder DAX.
    \item H2: Die Resistenz gegenüber Konflikten variiert stark zwischen verschiedenen regionalen Indizes (Geopolitische Resilienz).
    \item H3: Mithilfe von Clustering-Verfahren lassen sich Marktphasen identifizieren, die eindeutig mit Eskalationsstufen von Konflikten korrelieren.
\end{itemize}
Der wirtschaftliche Mehrwert liegt in einer verbesserten Risikobewertung für Investoren, um krisenbedingte Marktanomalien frühzeitig zu erkennen. (in Anlehnung an Google Gemini)

\subsection{Data Understanding}
Die Phase des Data Understanding befasst sich mit der Erhebung, Qualitätsprüfung und strukturellen Einordnung der Daten \parencite{chapman2000crisp}. Für die vorliegende Analyse werden tägliche Indexdaten mit ereignisbasierten Konfliktdaten auf Kalendertagebene zusammengeführt. Zentrale Prüfpunkte in dieser Phase sind die Datumsabdeckung beider Quellen, fehlende Handelswerte sowie die Präzision der Konfliktdatierung. Für die spätere Modellierung werden nur diejenigen Attribute berücksichtigt, die unmittelbar für das tagesgenaue Matching und die Konstruktion der finalen Analysemerkmale benötigt werden.

\subsubsection{Datenquellen und Umfang}
Der primäre Datensatz für die Finanzmarktanalyse stammt aus dem Kaggle-Stock-Exchange-Datensatz und enthält historische Tageswerte internationaler Indizes. Als zweite Quelle wird das UCDP GED (Global Event Dataset) verwendet, das konfliktbezogene Einzelereignisse mit Datums- und Intensitätsangaben bereitstellt. Für die tatsächliche Datenpipeline wurden die in Tabelle~\ref{tab:ucdp_used_attrs} und Tabelle~\ref{tab:index_used_attrs} aufgeführten Attribute genutzt.

\begin{table}[h!]
\centering
\caption{Genutzte Attribute aus den Konfliktdaten (UCDP GED)}
\label{tab:ucdp_used_attrs}
\begin{tabular}{@{}ll@{}}
\toprule
Spalte & Beschreibung \\ \midrule
date\_start & Ereignisdatum fuer das Matching auf Handelstage \\
date\_prec & Datumsqualität; verwendet als Filter (nur Wert 1) \\
type\_of\_violence & Konflikttyp zur täglichen Konfliktstruktur \\
country\_id & Anzahl betroffener Länder pro Tag \\
region & Anzahl betroffener Regionen pro Tag \\
event\_clarity & Qualitätsmerkmal des Ereignisses \\
where\_prec & Präzision der Ortsangabe \\
number\_of\_sources & Quellenanzahl als Robustheitsindikator \\
best & Zentrale Schätzung der Todesopfer (Intensität) \\
deaths\_civilians & Zivile Opfer pro Ereignis \\ \bottomrule
\end{tabular}
\caption*{\footnotesize Quelle: Eigene Darstellung}
\end{table}

\begin{table}[h!]
\centering
\caption{Genutzte Attribute aus den Indexdaten (Kaggle)}
\label{tab:index_used_attrs}
\begin{tabular}{@{}ll@{}}
\toprule
Spalte & Beschreibung \\ \midrule
Index & Eindeutige Indexkennung \\
Date & Handelstag und Matching-Key \\
Open & Eroeffnungskurs \\
High & Tageshoch \\
Low & Tagestief \\
Close & Schlusskurs \\
Adj\_Close & Bereinigter Schlusskurs \\
Volume & Handelsvolumen \\
daily\_return & Abgeleitete Tagesrendite \\
intraday\_range & Abgeleitete Tagesvolatilität \\
volume\_change & Abgeleitete Veränderung des Volumens \\ \bottomrule
\end{tabular}
\caption*{\footnotesize Quelle: Eigene Darstellung}
\end{table}

\subsection{Data Preparation}
Die Data-Preparation-Phase operationalisiert die im Data-Understanding identifizierten
Qualitaetsanforderungen in eine reproduzierbare Transformationspipeline.
Ziel ist die Erzeugung einer indexzentrierten Analysebasis mit exakt einer
Beobachtung pro Index und Handelstag. Diese Entscheidung ist methodisch zentral, da ein
naiver Ereignis-Join (ein Konfliktereignis auf mehrere Indizes desselben Tages) zu
Duplikaten, impliziter Uebergewichtung konfliktreicher Tage und damit zu verzerrten
Distanz- bzw. Regressionsstrukturen fuehren kann.

\subsubsection{Pipeline-Design}
Die Pipeline besteht aus vier aufeinander aufbauenden Schritten:
\begin{enumerate}
  \item \textbf{Aufbereitung der Indexdaten} \newline
  Die Rohdaten werden chronologisch je
  Index sortiert und um relative Marktmerkmale erweitert. Die zentrale Renditevariable
  wird als
  \[
  r_{i,t} = \frac{P_{i,t} - P_{i,t-1}}{P_{i,t-1}}
  \]
  definiert, wobei $P_{i,t}$ den Schlusskurs von Index $i$ am Handelstag $t$ bezeichnet.
  Ergaenzend werden \texttt{intraday\_range} als relative Tageshandelsspanne sowie
  \texttt{volume\_change} als relative Volumenaenderung berechnet.

  \item \textbf{Qualitaetsgesicherte Filterung der Konfliktdaten} \newline
  Aus \texttt{GEDEvent\_v25\_1.csv} werden nur fuer das Matching und die
  Konfliktcharakterisierung relevante Felder geladen. Anschliessend erfolgt ein strikter
  Filter auf \texttt{date\_prec = 1}. Dadurch werden nur Ereignisse mit hoher
  Datumsgenauigkeit in die Analyse aufgenommen; Messfehler durch unscharf datierte
  Ereignisse werden reduziert. Zudem werden technische Kodierungen fehlender Werte
  (z. B. \texttt{number\_of\_sources < 0}) in fehlende Werte ueberfuehrt,
  um verzerrte Mittelwerte zu vermeiden.

  \item \textbf{Tagesaggregation der Konfliktintensitaet} \newline
  Die Konfliktereignisse werden auf Tagesebene aggregiert
  (u. a. \texttt{conflict\_count}, \texttt{fatalities\_best\_sum},
  \texttt{countries\_affected}, \texttt{regions\_affected},
  \texttt{type\_of\_violence\_x\_count}). Damit wird aus einem ereignisbasierten in ein
  handelstagsbasiertes Merkmalsdesign ueberfuehrt. Diese Aggregation ist notwendig, um
  spaeter robuste Vergleichbarkeit zwischen ruhigen und eskalierenden Markttagen
  herzustellen.

  \item \textbf{Tagesgenaues Matching und fehlende Konflikte als Nullereignis} \newline
  Die aggregierten Konfliktmerkmale werden per Kalenderdatum auf die Indextabelle
  gemerged (Left-Join auf die Indexbasis). Tage ohne Konflikttreffer erhalten fuer
  Konfliktzaehl- und Intensitaetsvariablen den Wert 0 sowie die Indikatorvariable
  \texttt{has\_conflict = 0}. Dadurch bleibt die vollstaendige Handelstagsstruktur
  erhalten, was sowohl fuer Event-Study-Fenster als auch fuer ML-Modelle mit
  Klassifikationsziel essenziell ist.
\end{enumerate}

\subsubsection{Methodische Begruendung der Transformationen}
Die Wahl relativer statt absoluter Preismerkmale reduziert Skaleneffekte zwischen
heterogenen Indizes und verhindert, dass das Modell primaer nominale Kursniveaus lernt.
Der Praezisionsfilter in den GED-Daten stellt interne Validitaet sicher, weil nur
zeitlich hinreichend verortete Ereignisse als Exposition wirken. Die Tagesaggregation
erhoeht die statistische Stabilitaet, da sie Einzelereignisrauschen in robuste
Tagesindikatoren ueberfuehrt. Die Nullbelegung nicht betroffener Tage implementiert eine
konsistente Kontrollgruppe innerhalb derselben Zeitachse.

Die so erzeugte Datenbasis wird in den nachgelagerten Analysephasen unmittelbar
weiterverwendet, etwa fuer eventbasierte Renditeauswertungen sowie fuer
klassifikationsorientierte Modellierungsansaetze mit zusaetzlichen
modellspezifischen Vorbereitungsschritten (u. a. Skalierung,
Target-Konstruktion und Missing-Handling). Die Korrektheit der zentralen
Tagesaggregation wird durch formalisierte Plausibilitaets- und
Konsistenzpruefungen abgesichert, wodurch die Reproduzierbarkeit der
Data-Preparation nicht nur dokumentiert, sondern auch methodisch validiert wird.

\subsection{Modeling}
Einordnung der in Kapitel 3 beschriebenen Algorithmen in den Workflow.

\subsection{Evaluation}
Methodik der Modellgüteprüfung (Train-Test-Split, Metriken etc.).

\subsection{Deployment}
Beschreibung des Projekt-Repositorys und der Reproduzierbarkeit.

\section{Implementierung der Algorithmen}
\subsection{Lineare Regression}
\subsection{DBSCAN}
\subsection{KNN}
\subsection{K-Means Clustering}

Das K-Means-Verfahren wurde als exploratives, unueberwachtes Verfahren eingesetzt,
um ähnliche Markt-Konflikt-Situationen in der Datenbasis zu segmentieren. Ziel war
nicht die kausale Schätzung eines Effekts, sondern die Identifikation von
wiederkehrenden Strukturen im Merkmalsraum. Hierfür wurden numerische Merkmale wie
tägliche Rendite, Intraday-Range, Konfliktanzahl und Konfliktschwere verwendet.

Vor der Modellierung erfolgte eine Standardisierung aller Features, damit keine
einzelne Variable allein aufgrund ihrer Skala die Distanz dominiert. Anschließend
wurde K-Means für mehrere Werte von $k$ berechnet. Die Wahl der Clusterzahl
erfolgte datengetrieben über den Silhouette-Score. Formal minimiert K-Means die
quadratische Distanz der Beobachtungen zu ihren jeweiligen Clusterzentren:

\[
\min_{C_1,\ldots,C_k} \sum_{j=1}^{k} \sum_{x_i \in C_j} \lVert x_i - \mu_j \rVert^2
\]

mit $\mu_j$ als Zentrum von Cluster $C_j$.

Die Ergebnisse wurden über Clusterprofile (Mittelwerte je Cluster) interpretiert.
Dabei zeigte sich, dass K-Means insbesondere zwischen konfliktarmen und
konfliktreichen Handelstagen differenziert. Diese Erkenntnis ist für die
Hypothesengenerierung hilfreich, ersetzt jedoch keine inferenzstatistische
Wirkungsanalyse. Daher wird die kausale bzw. statistisch-signifikante Prüfung des
Zusammenhangs in der Arbeit über Event-Study- und Regressionsansätze ergänzt.

\section{Ergebnisse und Visualisierungen}
\section{Limitationen}
\section{Fazit}

\end{document}

